% Shared preamble for the three ENS003 software-team presentations.
% Each presentation \input{}'s this file.

\documentclass[11pt,aspectratio=169]{beamer}

\usepackage[utf8]{inputenc}
\usepackage[T1]{fontenc}
\usepackage{lmodern}
\usepackage{amsmath}
\usepackage{booktabs}
\usepackage{graphicx}
\usepackage{listings}
\usepackage{xcolor}
\usepackage{hyperref}

% Built-in Beamer theme so the deck compiles on any TeX Live / MacTeX
% install without external packages. Customise colours below.
\usetheme{Boadilla}
\usecolortheme{seahorse}
\setbeamertemplate{navigation symbols}{}
\setbeamertemplate{footline}[frame number]

\definecolor{ink}{HTML}{1A1A1A}
\definecolor{accentred}{HTML}{A02828}
\definecolor{accentgreen}{HTML}{2C662C}
\definecolor{codebg}{HTML}{F7F5EE}
\definecolor{rulegrey}{HTML}{888888}

\setbeamercolor{normal text}{fg=ink}
\setbeamercolor{frametitle}{bg=ink, fg=white}
\setbeamercolor{title}{fg=ink}
\setbeamercolor{author}{fg=ink}
\setbeamercolor{date}{fg=rulegrey}
\setbeamercolor{block title}{bg=ink, fg=white}
\setbeamercolor{block body}{bg=codebg, fg=ink}
\setbeamercolor{alerted text}{fg=accentred}
\setbeamercolor{structure}{fg=ink}
\setbeamerfont{frametitle}{series=\bfseries}

\lstset{
  basicstyle=\ttfamily\scriptsize,
  backgroundcolor=\color{codebg},
  language=Python,
  keywordstyle=\color{accentred},
  stringstyle=\color{accentgreen},
  commentstyle=\color{rulegrey}\itshape,
  breaklines=true,
  numbers=none,
  showstringspaces=false,
  frame=single,
  framerule=0pt,
  framesep=6pt,
  xleftmargin=0pt,
  xrightmargin=0pt
}

% Convenience commands
\newcommand{\code}[1]{\texttt{\small #1}}
\newcommand{\file}[1]{\textcolor{accentred}{\texttt{\small #1}}}
\newcommand{\unit}[1]{\,\mathrm{#1}}

% Title-slide metadata is set per-deck (so the same preamble can power
% both Turkish and English versions of each talk).


\newcommand{\courseHeader}{ENS003 --- Digital Twins for Health Sciences}
\newcommand{\courseSubheader}{Istinye University · Faculty of Engineering and Natural Sciences}
\newcommand{\instructor}{Instructor: Prof. Dr. Şenol Pişkin}
\newcommand{\projectTitle}{A Surrogate ML Model for a Ti-6Al-4V Cementless Hip Implant}

\title{Surrogate Model: GPR and 5-fold Grouped CV}
\subtitle{\projectTitle}
\author{Recep Kamil Fırat \\ \small Computer Engineering · 2309011079}
\date{May 2026 \\ \scriptsize \courseHeader{} · \instructor}

\begin{document}

\maketitle

% ---------------------------------------------------------------------------
\begin{frame}{In this talk}
  \begin{itemize}
    \item What the surrogate is for: pulling FEA output down to milliseconds
    \item Why Gaussian Process Regression (GPR), not XGBoost
    \item Cross-validation strategy: \textit{why} 5-fold grouped
    \item \(R^2_{\text{CV}}\) measured for all 10 targets
    \item Limits we report honestly
  \end{itemize}
  \vspace{0.4cm}
  \begin{block}{My responsibility (software team)}
    Surrogate model choice, hyperparameter tuning, cross-validation,
    \file{src/train.py} and \file{models/}.
  \end{block}
\end{frame}

% ---------------------------------------------------------------------------
\begin{frame}{The problem}
  \begin{columns}[T,onlytextwidth]
    \column{0.55\textwidth}
    The mechanical team solves each scenario with ANSYS Static Structural:
    \begin{itemize}\small
      \item One solve takes 5--10 minutes
      \item Covers a single mass/K/angle point
      \item \alert{Unsuitable} for an interactive dashboard
    \end{itemize}
    \vspace{0.4cm}
    Surrogate goals:
    \begin{itemize}\small
      \item \(\approx\) 1 ms response for the same inputs
      \item \alert{Reliable interpolation} inside the training grid
      \item Posterior \(\pm\sigma\) (calibrated uncertainty)
    \end{itemize}
    \column{0.45\textwidth}
    \begin{center}
      \footnotesize
      \begin{tabular}{lr}
        \toprule
        FEA solve & 5--10 min \\
        Surrogate predict & \(\approx\) 1 ms \\
        Speed-up & \(\sim 10^5\) \\
        \bottomrule
      \end{tabular}
    \end{center}
  \end{columns}
\end{frame}

% ---------------------------------------------------------------------------
\begin{frame}{Model selection --- 270 rows, target \(\mathrm{SF}_{\text{neck}}\)}
  Head-to-head benchmark on the same data, 5-fold grouped CV:
  \begin{center}\footnotesize
    \begin{tabular}{lrrl}
      \toprule
      Model & \(R^2_{\text{CV}}\) & MAE\(_{\text{CV}}\) & Unique outputs / 51 \\
      \midrule
      XGBoost (depth=1, stumps) & 0.986 & 0.032 & 18 \\
      XGBoost (depth=3)         & 0.988 & 0.024 & 27 \\
      Kernel Ridge (RBF)        & -0.005 & 0.369 & 5 \\
      \alert{\textbf{Gaussian Process (RBF + white)}} & \alert{\textbf{0.9999}} & \alert{\textbf{0.00077}} & \alert{\textbf{51}} \\
      \bottomrule
    \end{tabular}
  \end{center}

  \vspace{0.4cm}
  \textbf{Smoothness probe:} fix \(K=2.5, \theta=0°\), sweep the mass
  slider over 51 points. GPR produces a distinct value at every point
  (continuous curve); tree-based methods collapse onto few distinct
  values \(\Rightarrow\) \alert{stepped} function.

  \vspace{0.3cm}
  \footnotesize
  On 162 unique grid points, XGBoost averages \(\sim 6\) samples per leaf,
  which is exactly what produces the step-function behaviour.
\end{frame}

% ---------------------------------------------------------------------------
\begin{frame}[fragile]{GPR architecture}
\begin{lstlisting}
from sklearn.gaussian_process import GaussianProcessRegressor
from sklearn.gaussian_process.kernels import RBF, ConstantKernel, WhiteKernel

kernel = (
    ConstantKernel(1.0, (1e-3, 1e3))
    * RBF(length_scale=1.0, length_scale_bounds=(1e-2, 1e3))
    + WhiteKernel(noise_level=1e-3, noise_level_bounds=(1e-6, 1e-1))
)

base = GaussianProcessRegressor(
    kernel=kernel, normalize_y=True,
    n_restarts_optimizer=5, random_state=42, alpha=1e-8,
)

model = TransformedTargetRegressor(
    regressor=base, transformer=StandardScaler(),
)
\end{lstlisting}

  \begin{itemize}\small
    \item Kernel hyperparameters fit by marginal-likelihood maximisation
    \item \code{normalize\_y} \(+\) \code{TransformedTargetRegressor} \(\Rightarrow\) target-scale invariant
    \item Closed-form posterior \(\Rightarrow\) every prediction comes with \(\sigma\) for free
  \end{itemize}
\end{frame}

% ---------------------------------------------------------------------------
\begin{frame}{CV strategy: \textit{why} grouped?}
  \textbf{Problem:} 108 of 270 rows are ANSYS re-runs --- the same
  \((m, K, \theta)\) point repeated with tiny numerical differences.

  \vspace{0.3cm}
  \begin{itemize}
    \item With plain \code{KFold}\,/\,\code{LOO}, the same grid point can
          land in both train and test folds \(\Rightarrow\) \alert{leakage}
    \item Visible symptom: \(R^2\) inflates artificially to \(\approx 1.0\)
    \item We lose the ability to measure honest out-of-sample performance
  \end{itemize}

  \vspace{0.4cm}
  \textbf{Fix:} \code{GroupKFold(n\_splits=5)}, with the group ID coming
  from Ata's \code{grid\_groups()} helper:
  \begin{center}
    \texttt{"\,m\,|\,K\,|\,theta\,"} \quad \(\Rightarrow\) \quad re-runs share a fold
  \end{center}
\end{frame}

% ---------------------------------------------------------------------------
\begin{frame}[fragile]{The training loop}
\begin{lstlisting}
from sklearn.model_selection import GroupKFold

def _grouped_cv_eval(X, y, groups, n_splits):
    preds = np.full(len(y), np.nan)
    for tr, te in GroupKFold(n_splits).split(X, y, groups=groups):
        model, _ = _make_estimator()
        model.fit(X.iloc[tr], y.iloc[tr])
        preds[te] = model.predict(X.iloc[te])
    return {
        "r2_cv": r2_score(y, preds),
        "mae_cv": mean_absolute_error(y, preds),
    }
\end{lstlisting}

  \vspace{0.3cm}
  Per target:
  \begin{enumerate}\small
    \item Out-of-fold predictions across 5 folds
    \item Final fit on the full 270 rows \(\Rightarrow\) \file{models/<target>.joblib}
    \item Metrics written to \file{models/metrics.json}
  \end{enumerate}
\end{frame}

% ---------------------------------------------------------------------------
\begin{frame}{CV results --- 10 targets}
  \begin{center}\footnotesize
    \begin{tabular}{lrr}
      \toprule
      Target & \(R^2_{\text{CV}}\) & MAE\(_{\text{CV}}\) \\
      \midrule
      safety\_factor\_neck\_min            & 0.9999 & 7.7e-4 \\
      safety\_factor\_stem1\_min           & 0.9999 & 8.3e-4 \\
      safety\_factor\_stem2\_min           & 0.9997 & 4.4e-3 \\
      max\_equivalent\_vonmises\_stress\_Pa & 0.9998 & 0.53 MPa \\
      avg\_equivalent\_vonmises\_stress\_Pa & 0.9998 & 10.6 kPa \\
      \alert{min\_equivalent\_vonmises\_stress\_Pa} & \alert{0.034} & \alert{(FEA noise floor)} \\
      max\_principal\_stress\_Pa            & 0.9998 & 0.57 MPa \\
      avg\_principal\_stress\_Pa            & 0.9998 & 5.0 kPa \\
      min\_principal\_stress\_Pa            & 0.9998 & 0.15 MPa \\
      max\_total\_deformation\_m            & 0.9998 & \(7.4 \times 10^{-7}\) m \\
      \bottomrule
    \end{tabular}
  \end{center}

  \vspace{0.4cm}
  9/10 targets clear \(R^2 > 0.999\); \code{min\_equivalent} sits at
  \(\sim 10^{-10}\) --- the FEA solver's noise floor. We report it honestly
  and hide it from the dashboard.
\end{frame}

% ---------------------------------------------------------------------------
\begin{frame}{Honest limits}
  \begin{itemize}
    \item \textbf{Interpolation is reliable, extrapolation is not.}
          Predictions for \(m < 50\unit{kg}\), \(K > 6\), \(\theta > 20°\)
          should be treated with care.
    \item \textbf{18-face fixed support is constant.} Different fixation
          patterns are not modelled; new FEA grid would be needed.
    \item \textbf{\(F_y = 0\) in every row.} No out-of-plane loading;
          real activities are mostly in-plane so this is acceptable,
          but not complete.
    \item \textbf{Kernel hyperparameters near the bounds for some targets
          (ConvergenceWarning).} Marginal impact on \(R^2\); the bounds
          can be widened later.
  \end{itemize}
  \vspace{0.4cm}
  \footnotesize
  Fallback chain: \code{GPR \(\to\) XGBoost \(\to\) sklearn-GBR} --- if
  sklearn's GPR is unavailable in an environment, the others kick in
  (\code{test\_make\_estimator\_falls\_back\_to\_sklearn\_gbr}).
\end{frame}

% ---------------------------------------------------------------------------
\begin{frame}[fragile]{Demo}
\begin{lstlisting}[language=bash, basicstyle=\ttfamily\small]
$ ./run.sh train
[...]
$ python -c "
import json
m = json.loads(open('models/metrics.json').read())
for r in m:
    print(f'{r[\"target\"]:42s}  R2_cv={r[\"r2_cv\"]:.4f}')
"
safety_factor_neck_min                       R2_cv=0.9999
safety_factor_stem1_min                      R2_cv=0.9999
safety_factor_stem2_min                      R2_cv=0.9997
[...]
\end{lstlisting}
  \vspace{0.4cm}
  Test suite: \alert{27 passed}, 96\% coverage. The dashboard reads
  \file{models/} and runs all 10 models per predict call.
\end{frame}

% ---------------------------------------------------------------------------
\begin{frame}{Questions}
  \vfill
  \begin{center}
    \Large Thank you.\\[0.4cm]
    \normalsize Questions?
  \end{center}
  \vfill
  \scriptsize
  \begin{itemize}
    \item \file{src/train.py} (GroupKFold + GPR + XGB fallback)
    \item \file{models/metrics.json} (live CV scores)
    \item \file{tests/test\_train\_predict.py}
    \item Development log: \file{docs/PROJECT\_NOTES.md} Sprints 3--4
  \end{itemize}
\end{frame}

\end{document}
